\chapter{Introduzione}

L'informatica teorica si occupa di due tipi di oggetti profondamente diversi tra loro. Da un lato i \textbf{dati}: strutture finite, costruite a partire da componenti più semplici in un numero finito di passi. Dall'altro ci sono i \textbf{processi}: sistemi che evolvono nel tempo, che possono non terminare mai, e la cui natura sta in come si comportano quando vengono osservati.

La matematica classica offre uno strumento potente per ragionare sul primo tipo di oggetti: l'\textbf{induzione}. La struttura induttiva permette sia di \emph{definire} funzioni per ricorsione sulla struttura dei dati, sia di \emph{dimostrare} proprietà per induzione.

I processi, però, non si lasciano descrivere altrettanto bene da questo stesso strumento: questi oggetti non si costruiscono a partire da casi base, perché non hanno un caso base. Ciò che li caratterizza è \emph{cosa succede quando li si osserva}: quale simbolo producono, in quale stato transitano, se si arrestano o continuano. Per ragionare su oggetti di questo tipo serve un principio duale all'induzione, la \textbf{coinduzione}, che permetta di definire funzioni per \emph{coricorsione}, specificandone il comportamento passo dopo passo anziché costruendole dal basso, e di dimostrare l'uguaglianza di due comportamenti mostrando che sono osservativamente indistinguibili, invece di scomporli nei loro costituenti.

La \textbf{teoria delle categorie} fornisce il linguaggio in cui questa dualità tra induzione e coinduzione si può esprimere in modo del tutto preciso e generale. Un tipo di dato induttivo è catturato dalla nozione di \textbf{algebra iniziale} per un funtore: la struttura più semplice che si può costruire a partire da un insieme di costruttori. Un sistema di transizione, dualmente, è catturato dalla nozione di \textbf{coalgebra}, e il suo comportamento osservabile è descritto dalla \textbf{coalgebra finale}: la struttura più informativa verso cui ogni altro sistema dello stesso tipo si proietta in modo unico.

Questa prova finale si propone di sviluppare, in modo autocontenuto, gli strumenti categorici necessari a formalizzare questo punto di vista, e di applicarli a tre esempi classici dell'informatica teorica. Il capitolo 2 introduce i concetti di base della teoria delle categorie — categorie, funtori, isomorfismi, oggetti finali — necessari a dare un significato preciso alla nozione di coalgebra. Il capitolo 3 sviluppa la teoria generale delle coalgebre: la definizione di coalgebra e dei suoi omomorfismi, la coalgebra finale come oggetto universale, i principi di coricorsione e coinduzione che essa induce e la nozione di bisimulazione come criterio di equivalenza comportamentale. Il capitolo 4, infine, applica questi strumenti a tre ambiti concreti dell'informatica teorica — i linguaggi regolari, le grammatiche libere dal contesto e le macchine di Turing — mostrando come proprietà classiche di ciascuno di questi formalismi (chiusura dei linguaggi regolari, raggiungibilità e produttività dei simboli di una grammatica, correttezza di simulazioni tra varianti di macchine di Turing) si lascino esprimere e dimostrare naturalmente nel linguaggio delle coalgebre, degli operatori temporali coinduttivi e della bisimulazione.
